\documentclass{beamer}

\usepackage{bm}
\usepackage{amsmath}
\usepackage{amssymb}
\usepackage{xcolor}
\usepackage{graphicx}
%\usepackage{tikz}

\title{Linear Transformations}
\subtitle{\rotatebox{-27}{They} \scalebox{2}[3]{can} be \rotatebox{27}{great}}
\author{Ahmed Mohammed}
\institute{Georgia Institute of Technology}
\date{PLUS Session, September 2023}

\usetheme{Berlin}
\usefonttheme{serif}
\usecolortheme{seagull}

\AtBeginSection[] {
	\begin{frame}
		\frametitle{Table of Contents}
		\tableofcontents
	\end{frame}
}

\DeclareMathOperator{\vspan}{Span}
\renewcommand{\vec}[1]{\bm{#1}}

\begin{document}
	\begin{frame}
		\maketitle
	\end{frame}
	
	\section{Open Activity}
	\begin{frame}
		\frametitle{Revision: \textit{What is\dots}}
		\begin{enumerate}[<+(1)->]
			\item What is a \textit{linear} transformation from \(\mathbb{R}^n\) to \(\mathbb{R}^m\)?
			\item What is the difference between linear transformations and other transformations from \(\mathbb{R}^n\) to \(\mathbb{R}^m\)?
			\item What do we call the whole set of inputs of a transformation \(T \colon \mathbb{R}^n \to \mathbb{R}^m\)?
			\item What is \(\mathbb{R}^m\) in \(T \colon \mathbb{R}^n \to \mathbb{R}^m\) called?
			\item What is the difference between the \emph{range} of a transformation \(T\colon\mathbb{R}^n\to\mathbb{R}^m\) and its \emph{codomain}?
		\end{enumerate}
	\end{frame}

	\section{Main Activity}
	\begin{frame}
		\frametitle{(Main?) Activity}
		\centering Look at the worksheet!
	\end{frame}

	\section{Closing Activity}
	\begin{frame}
		\frametitle{Looking at the \(\vec{e}_i\)}
		\begin{itemize}
			\item<1-> We've talked about how linear transformations look, how they behave, and what their components are.
			\item<2-> Next, we see how we can represent \textit{linear transformations from \(\mathbb{R}^n \to \mathbb{R}^m\)}.
		\end{itemize}
	\end{frame}
\end{document}